\cotefile{37}
\scanpage{52}
\noindent SGA 5 \(\overline{\underline{\mathrm{VII}}}\)\quad th de Lefschetz faible sans supposer \(X\) ni \(Y\) lisses (seulement \(X-Y\) lisse \ldots)

\noindent SGA 5 \(\overline{\underline{\unsure{\mathrm{V}|\mathrm{VI}}}}\)\quad \(\mathbb{Q}_{\ell}\)-faisceau : rajouter définition raisonnable en termes de \uline{champ} \uline{associé}

\noindent SGA 5. \(\deleted{\overline{\underline{\unsure{\mathrm{IV}|\mathrm{V}|\mathrm{VI}}}}}\)\,\inserted{\tikz[baseline=(vi.base)]{\node[inner sep=1pt] (vi) {\(\overline{\underline{\mathrm{VI}}}\)};\draw[thin] ([xshift=-1pt,yshift=1pt]vi.north west) .. controls +(270:0.6cm) and +(265:0.6cm) .. ([xshift=3pt,yshift=3pt]vi.north east);}}\quad Reprendre en langage de \(\mathbb{Z}_{\ell}\)-faisceaux ou\textsuperscript{\inserted{\unsure{*|+}}} \(\mathbb{Q}_{\ell}\)-faisceaux les notions essentielles dont il a été question dans SGA 4 (, Ajouter nouveau N\textsuperscript{\uline{o}}).

\noindent SGA 5 \(\overline{\underline{\mathrm{IV}}}\)\quad Essayer de faire dans le cas relativement lisse ttes les constructions essentielles

\figure{\begin{tikzpicture}[x=\linewidth,y=1cm]
\draw[thin] (0,0) .. controls (.23,.12) and (.49,.08) .. (.62,.06)
 .. controls (.77,.05) and (.90,-.03) .. (.96,.02);
\end{tikzpicture}}

\noindent \(\overline{\underline{\mathrm{VI}}}\)\quad Systèmes projectifs \inserted{\tikz[baseline=(adiques.base)]{\node[inner sep=0pt] (adiques) {adiques};\draw[thin] ([xshift=-0.35em,yshift=-0.7ex]adiques.south west) .. controls +(85:1.2ex) and +(-1em,-0.15ex) .. (adiques.south) .. controls +(0.8em,0.15ex) and +(0em,-0.7ex) .. ([xshift=0.12em,yshift=0.1ex]adiques.east);}} dans les catégories dérivées \ldots\ (ou dans thèse)\margin{\unsure{:}}
% Two small isolated dots slant upward to the right in the right margin.

\noindent SGA 5 \(\mathrm{V}\)\quad dans introduction du \unsure{livre|Hom} ou de l'expo\unsure{sé} signaler l'ennui des coeff. trop particuliers \((\mathbb{Z}_{\ell},\,\mathbb{Z}/\ell^{\unsure{N|n|\nu}}\mathbb{Z}\ -\,-)\)

\figure{\begin{tikzpicture}[x=\linewidth,y=1cm]
\draw[thin] (0,.03) .. controls (.12,.11) and (.19,.04) .. (.29,-.02)
 .. controls (.39,-.08) and (.45,.02) .. (.57,.01)
 .. controls (.74,.02) and (.85,-.11) .. (.91,-.06);
\end{tikzpicture}}

\noindent \(\mathrm{I}\)\quad \S\ 5\qquad Donner, si \(i\unsure{{<}|{:}|{\in}}\eta\overset{\unsure{\sim}}{\longrightarrow} X\) inclusion du \(\mathrm{pt}\) générique, et \(F_{\eta}\) Module loc. libre sur \(\eta\)
% A faint bowed stroke above the arrow shaft is retained as an uncertain tilde; its intentional status is unresolved.
\[
 D\bigl(i_{*}(F_{\eta})\bigr)=i_{*}\bigl(D(F_{\eta})\bigr)
\]
% archive sheet number: 51
\clearpage
